% 第24段階の追加記録. 第23段階の本文に続けて追加する.
% 既存の \bm, \bbo, proposition, proof, booktabs, hyperrefを使用する.
% 新しい引用キー, 文献環境, 文書開始・終了命令は追加しない.
% 根拠資料: bivariate_survival_reconstruction_stage24.zip.
% 第25段階の大域最適値の検証とは分けて記録する.

\section{第24段階の目的と再解析する親クラス}
\label{sec:stage24_scope}

第23段階で共通尾部クラスを外れた13件を全て再解析する.
粗い記録・深度3の5件と深度4の8件であり,
同じ母標本の再観測を含むため, 母標本の種類は8個である.
新しい親標本を生成せず, 観測区間, 限界, 学習・検証個数,
全実数空間の観測同値類と共通初期測度$R$を保持する.
全データ個数を点推定に用いるが,
学習側だけの2EM予測, 検証尤度閾値と親信頼集合は変更しない.

原EMの13結果は非負・規格化された正当な生存分布である.
今回の問題は, 別途宣言した尾部クラスへの不所属であり,
分布自体の不正ではない. 旧結果を上書きしない.
各入力の観測同値類は16個であり, 親クラスを
\begin{align}
\mathcal K
&=\{\bm m\ge0:\bm1^{\mathsf T}\bm m=1,
\ \bm D_j\bm m\le1/2\quad(j=1,2,3)\}
\label{eq:stage24_parent_class}
\end{align}
とする. $\bm D_j$はD領域への所属指示行である.
対象設計では$D_3\subset D_1\cap D_2$なので, 第三制約は冗長である.
この尾部上界は外部のモデル条件であり, 観測から推定した量ではない.

\section{同じ補完E-stepと制約付きM-step}
\label{sec:stage24_estep}

正の個数$n_i$を持つ記録について,
$a_i(\bm m)=\bm A_i\bm m$, $q_i(\bm m)=\bm V_i\bm m$とする.
従来のD補完で得られる期待潜在個数と正規化は,
\begin{align}
\nu_h(\bm m)
&=m_h\sum_i n_i\left[
\frac{A_{ih}}{a_i(\bm m)}+
\frac{1-V_{ih}}{q_i(\bm m)}\right],
\label{eq:stage24_expected_counts}\\
K(\bm m)&=\sum_i\frac{n_i}{q_i(\bm m)}
=\sum_h\nu_h(\bm m),\qquad
w_h=\frac{\nu_h(\bm m)}{K(\bm m)}.
\label{eq:stage24_normalization}
\end{align}
ここまでは元のEMと同じである.
制約付きM-stepでは, 正規化したsurrogate
\begin{align}
\overline{\mathcal Q}(\bm u\mid\bm m)
&=\sum_h w_h\log u_h
\label{eq:stage24_surrogate}
\end{align}
を$\mathcal K$上で最大化する.
$w_h=0$の項は加えない. 元の条件付き尤度が非凹でも,
このsurrogateは$\bm u$について凹である.
観測機構と補完E-stepは変更していないため,
可行な観測許容質量$\bm m,\bm u$について,
\begin{align}
\ell(\bm u)-\ell(\bm m)
&\ge K(\bm m)\{
\overline{\mathcal Q}(\bm u\mid\bm m)
-\overline{\mathcal Q}(\bm m\mid\bm m)\}
\label{eq:stage24_em_lower_bound}
\end{align}
という同じEM下界を使える.

\section{二つのD所属による四群への正確な還元}
\label{sec:stage24_four_groups}

$\mathcal H_{ab}=\{h:(D_{1h},D_{2h})=(a,b)\}$を用い,
\begin{align}
W_{ab}&=\sum_{h\in\mathcal H_{ab}}w_h,\qquad
p_{ab}=\sum_{h\in\mathcal H_{ab}}u_h
\label{eq:stage24_group_masses}
\end{align}
とする. 本段階では四つの$W_{ab}$が全て正の場合を扱う.
群総質量$p_{ab}$を固定したとき, 群内部の最適配分は
\begin{align}
u_h^*&=p_{ab}\frac{w_h}{W_{ab}}
\quad(h\in\mathcal H_{ab})
\label{eq:stage24_within_group}
\end{align}
である. 実際, 群内の条件付き質量を$r_h=u_h/p_{ab}$と書けば,
その項は$W_{ab}\log p_{ab}+\sum_{h\in\mathcal H_{ab}}w_h\log r_h$であり,
後者は$r_h=w_h/W_{ab}$で最大となる.
従って16質量のM-stepは,
\begin{align}
\max_{\bm p}\quad&\sum_{a,b}W_{ab}\log p_{ab},
\label{eq:stage24_group_objective}\\
p_{ab}&\ge0,\qquad \sum_{a,b}p_{ab}=1,\notag\\
p_{10}+p_{11}&\le\frac12,\qquad
p_{01}+p_{11}\le\frac12
\label{eq:stage24_group_constraints}
\end{align}
へ正確に還元される.
D所属の二成分の独立性を仮定したのではない.
元のA, B, C記録と補完期待量は$W_{ab}$に保持される.

\subsection{四つの活性集合に対応する候補}
\label{subsec:stage24_active_sets}

制約が非活性なら$p_{ab}=W_{ab}$である.
第一制約だけが活性の場合は$t=W_{10}+W_{11}$として,
\begin{align}
p_{1b}&=\frac{W_{1b}}{2t},\qquad
p_{0b}=\frac{W_{0b}}{2(1-t)}
\label{eq:stage24_first_active}
\end{align}
を得る. 第二制約だけの場合は両座標を交換する.
両制約が活性の場合は,
\begin{align}
p_{00}=p_{11}&=\frac{W_{00}+W_{11}}2,\qquad
p_{01}=p_{10}=\frac{W_{01}+W_{10}}2
\label{eq:stage24_both_active}
\end{align}
である.

各候補について可行性に加え, 非負な乗数$\beta_1,\beta_2$と$\eta$による
\begin{align}
\frac{W_{ab}}{p_{ab}}
&=\eta+\beta_1a+\beta_2b,
\label{eq:stage24_group_stationarity}\\
\beta_1(p_{10}+p_{11}-1/2)&=0,\qquad
\beta_2(p_{01}+p_{11}-1/2)=0
\label{eq:stage24_group_complementarity}
\end{align}
を有理数で検査する.

\begin{proposition}[四群M-stepの大域最大化]
\label{prop:stage24_mstep_optimality}
上の可行性, 非負乗数と相補性を満たす候補は,
四群問題および元の制約付きsurrogateの大域最大化解である.
\end{proposition}
\begin{proof}
凹関数の接線上界から, 任意の可行な$\bm p'$について,
\begin{align}
\sum_{a,b}W_{ab}\log\frac{p'_{ab}}{p_{ab}}
&\le\sum_{a,b}\frac{W_{ab}}{p_{ab}}(p'_{ab}-p_{ab})\notag\\
&=\beta_1\{p'_{10}+p'_{11}-p_{10}-p_{11}\}\notag\\
&\quad+\beta_2\{p'_{01}+p'_{11}-p_{01}-p_{11}\}\notag\\
&\le0.
\label{eq:stage24_mstep_proof}
\end{align}
正規化で$\eta$の項は消え, 最後の不等式は相補性と可行性による.
群内部の最適配分を戻せば元のsurrogateについても成立する.
\end{proof}

零の群期待質量を含む退化ケースや, 異なる尾部上界に対する一般解を,
今回の閉形式実装で扱ったとはしない.
対象13件の全1040更新では, 四群の正値性を確認した.

\section{有限精度GEMの採用条件と更新誤差}
\label{sec:stage24_finite_precision}

有理数反復の分母増大を抑えるため,
各回の厳密なM-step解を22有効十進桁の有理数候補へ変換する.
規格化後, 丸めによる尾部違反がある場合に限り,
固定した可行初期質量との有理数混合で新候補を作る.
これは採用前の候補構成であり,
旧クラス外推定値を事後的に修復して報告する操作ではない.
親モデルに正の質量floorを追加しない.

候補$\bm u$の採用には, 元の全制約と,
\begin{align}
\overline{\mathcal Q}(\bm u\mid\bm m)
-\overline{\mathcal Q}(\bm m\mid\bm m)&\ge0,\notag\\
\frac{\ell(\bm u)-\ell(\bm m)}{N}
&\ge\frac{K(\bm m)}N
\{\overline{\mathcal Q}(\bm u\mid\bm m)
-\overline{\mathcal Q}(\bm m\mid\bm m)\}\ge0
\label{eq:stage24_acceptance}
\end{align}
を有理数対数区間で検査する.
ここで$N=\sum_i n_i$は実際の観測総数であり, 補完数$K$とは異なる.
実装した反復は, 近似候補の増加を証明したGeneralized EMである.
厳密なM-step解$\bm u^*$との目的値差にも,
\begin{align}
0&\le\overline{\mathcal Q}(\bm u^*\mid\bm m)
-\overline{\mathcal Q}(\bm u\mid\bm m)\notag\\
&\le\sum_{h:w_h>0}w_h\left(\frac{u_h^*}{u_h}-1\right)
\label{eq:stage24_mstep_error}
\end{align}
という接線上界を用いる.

13件それぞれで固定80回, 合計1040更新を検査した.
全更新で可行性, surrogate増加と観測尤度増加を確認し,
一更新の平均対数尤度増加の最小下界は約$1.023829\times10^{-7}$,
正規化M-step損失の最大上界は約$2.452849\times10^{-24}$だった.
後者はM-step内部の計算誤差であり, 統計的な推定精度ではない.
80更新時点の制約付き一階残差上界は
約$5.68\times10^{-5}$から$5.33\times10^{-4}$であり,
この固定回数の結果を収束済みとはしない.

\section{別経路の境界候補と制約付き停止診断}
\label{sec:stage24_stationarity}

GEM監査とは別に, 同じ観測尤度を三つの決定論的初期値からSLSQPで最大化し,
活性面上のNewton補正を候補提案に使った.
初期値は80更新後, 共通$R$, 16クラス一様である.
提案を非負有理数へ変換し, 元の尾部制約,
正個数記録の正値性と80更新後からの観測尤度増加を検査した.
この部分をGEMの一更新とは呼ばない.

平均対数尤度の勾配を$\bm g(\bm m)=\nabla\ell(\bm m)/N$とする.
尤度の零次同次性により$\bm g^{\mathsf T}\bm m=0$であり,
制約付きの最大一次変化を
\begin{align}
\delta(\bm m)
&=\max_{\bm v\in\mathcal K}
\bm g(\bm m)^{\mathsf T}(\bm v-\bm m)
\label{eq:stage24_directional_gap}
\end{align}
と定義する.
非負な$\lambda_1,\lambda_2$と$z$が全質量列で
\begin{align}
\lambda_1D_{1h}+\lambda_2D_{2h}+z&\ge g_h
\label{eq:stage24_directional_dual}
\end{align}
を満たせば,
\begin{align}
0\le\delta(\bm m)&\le z+\frac{\lambda_1+\lambda_2}{2}
\label{eq:stage24_directional_upper}
\end{align}
となる. 数値LPが提案した乗数を非負有理数にした後,
$z=\max_h(g_h-\lambda_1D_{1h}-\lambda_2D_{2h})$と置き直し,
全列で双対可行性を検査する.

\begin{table}[htbp]
\centering
\caption{第24段階で別経路から得た制約付き候補の停止診断.
一階残差は平均対数尤度に対応する.}
\label{tab:stage24_candidates}
\begin{tabular}{crcc}
\toprule
深度 & 件数 & 旧候補の最小収録確率 & 新候補の最大一階残差上界 \\
\midrule
3 & 5 & 約$0.45464$--$0.48946$ & $9.45\times10^{-15}$ \\
4 & 8 & 約$0.42635$--$0.49094$ & $5.85\times10^{-14}$ \\
\bottomrule
\end{tabular}
\end{table}

新しい13候補は全て$q_j\ge1/2$を満たす一方,
同じ候補の未制約平均正スコアには最大約0.032846が残った.
クラス外へ出る方向も数える未制約スコアを,
制約付き問題の停止判定へ流用してはいけない.
また, 元の尤度は一般に非凹なので,
$\delta$を大域最適値との差と読み替えない.
本段階で証明した大域最大化は各surrogateについてであり,
元の尤度に対する13個の大域MLEを認定したものではない.

\section{不可視質量の役割と比較の範囲}
\label{sec:stage24_invisible_mass}

旧候補から完全不可視質量$\theta$を除いて規格化すれば,
条件付き記録比率は変わらず, 収録確率は$q_j/(1-\theta)$となる.
しかし13件のうち, この操作だけで尾部クラスへ入るのは1件だった.
残る12件では不可視質量を全て除いても最小収録確率が$1/2$未満となる.
違反を一律に初期不可視質量だけの問題とは解釈できない.

制約付きM-stepでは完全不可視セルにも,
\begin{align}
u_h^*&=\frac{m_h}{\eta+\beta_1+\beta_2}
\label{eq:stage24_invisible_update}
\end{align}
が適用され得る.
非活性の場合は元の保存則へ戻るが, 活性制約があれば一般には保存されない.
これは追加親クラス内で最大化する結果であり, 新しいD情報の観測ではない.
候補の不可視質量がゼロ付近になっても, 一意な識別を結論しない.

参考生存RMSEは13件全てで旧候補より小さくなったが,
クラス違反例だけを選んだ再解析である.
一般的な平均リスク改善や被覆率改善とは解釈しない.
旧未制約候補より尤度が下がる場合も, 実行可能集合が異なるので矛盾しない.
尤度の単調性を比較するのは, 同じ制約クラス内の可行更新間である.

\section{独立検証と保存範囲}
\label{sec:stage24_record}

独立検証器は公開入力から13モデルを再構成し,
元の学習2EM予測との一致, 1040更新,
16640個のM-step KKT座標等式と13境界候補を確認した.
M-step提案関数を呼ばず, 保存乗数から別経路で解を再構成する.
LP, 非線型最適化と乱数生成を呼ばずに検査し,
五種類の有害な改ざんを拒否した.
四群の256有理数対照では, 非活性90, 第一のみ57,
第二のみ57, 両方活性52件を確認した.
別のSLSQPとの16比較における目的差は最大$6.67\times10^{-16}$未満だった.

保存時には全13ケースを別出力先で再実行し,
13証拠JSONと比較CSVの全バイト一致,
集約JSONの実行時間以外の一致を確認している.
本LaTeX記録の作成時にも保存された13証拠を独立検証器で再検査した.
今回の再確認ではGEM反復と境界候補の最適化を再実行していない.

保存パッケージは
\nolinkurl{bivariate_survival_reconstruction_stage24.zip}である.
\texttt{core.py}に提案計算と正確なM-step,
\texttt{audit\_stage24.py}に監査,
\texttt{verify\_stage24.py}に独立検証を保存した.
\begin{verbatim}
python verify_stage24.py --outdir outputs
\end{verbatim}
比較表は\nolinkurl{outputs/comparison.csv},
証拠は\nolinkurl{outputs/certificates/}にある.
新しい母標本生成, 720件全体の再適合,
親信頼集合の変更と被覆率の再評価は行っていない.
次段階では, 一階条件を満たした同じ13候補について,
元の条件付き尤度の大域最適値との差を上下界によって検査する.
